% Vietnamese translation for OI Public Library
% Source-File: IOI中国国家候选队论文/2005/钱自强/钱自强.ppt
\documentclass[11pt]{article}
\usepackage{oipl-vn}

\title{Bản trình chiếu: Thuật toán di truyền}
\author{Qian Ziqiang}
\date{2005}

\begin{document}
\maketitle

\origtitle{Slide companion for genetic algorithms}
\sourcepath{IOI China candidate team papers / 2005 / Qian Ziqiang / Qian Ziqiang.ppt}

\section{Phạm vi bản dịch}

Bản trình chiếu đi cùng bài viết về thuật toán di truyền. Phần chữ đọc được
nhắc tới John Holland, \textit{Genetic Algorithm}, IOI 2005, cùng các thuật
toán Prim và Kruskal trong phần ví dụ.

\section{Tư tưởng}

Thuật toán di truyền mô phỏng tiến hóa: mỗi cá thể biểu diễn một nghiệm, hàm
thích nghi đánh giá chất lượng, rồi quần thể được cập nhật qua chọn lọc, lai
ghép và đột biến. Cách tiếp cận này thường dùng cho bài tối ưu khó, nơi lời
giải chính xác quá đắt hoặc không cần thiết.

\section{Ví dụ cây khung}

Việc nhắc tới Prim và Kruskal cho thấy slide có ví dụ liên quan tới cây khung
hoặc tối ưu trên đồ thị. Trong bối cảnh thuật toán di truyền, một cá thể có
thể mã hóa một cây hoặc một thứ tự chọn cạnh; hàm thích nghi đo tổng trọng số
và mức vi phạm ràng buộc.

\section{Kết luận}

Thuật toán di truyền mạnh ở khả năng tìm nghiệm tốt trong không gian lớn,
nhưng cần thiết kế mã hóa và hàm thích nghi cẩn thận. Bài trình bày nhấn mạnh
đây là công cụ heuristic, không phải thay thế cho chứng minh tối ưu khi bài
đòi hỏi đáp án chính xác.

\end{document}
